\documentclass[10pt,a4paper]{amsart}
\usepackage[margin=19mm]{geometry}
\usepackage{amsmath,amssymb,mathtools}
\usepackage[scr=rsfs,cal=euler]{mathalfa}
\usepackage{xfrac}
\usepackage[unicode,hidelinks,bookmarks=false]{hyperref}
\newcommand{\ie}{i.e.\ }
\newcommand{\cf}{cf.\ }
\newcommand{\resp}{respectively\ }
\newcommand{\Subsection}{\subsection}
\providecommand{\nobreakdash}{\nobreak\hbox{-}}
\setlength{\emergencystretch}{2em}
\multlinegap0pt
\usepackage{bm}
\usepackage{bbm}
\usepackage{dsfont}
\usepackage{xspace}
\usepackage{cancel}
\usepackage{mathtools}
\usepackage{amsthm}
\makeatletter
\@ifundefined{proposition}{\newtheorem{proposition}{Proposition}}{}
\@ifundefined{theorem}{\newtheorem{theorem}{Theorem}}{}
\makeatother
\title[The Matching Ramsey Number of Hypergraphs, Revisited]{The Matching Ramsey Number of Hypergraphs, Revisited}
\author{Saeed Shaebani, Meysam Alishahi}
\date{Source: arXiv:2101.04701v2 (CC BY 4.0)}
\begin{document}
\maketitle

\noindent\emph{Source and licence.} The Matching Ramsey Number of Hypergraphs, Revisited. Version arXiv:2101.04701v2 (CC BY 4.0), distributed under the Creative Commons Attribution 4.0 International licence, as recorded in the arXiv licence metadata for this exact version and shown on the abstract page https://arxiv.org/abs/2101.04701v2. Full rights notice: licenses/arxiv-2101.04701-CC-BY-4.0.txt.\par\smallskip

\noindent\textit{Editorial scope.} Counted unit: the Proposition whose source label is \texttt{None} in the article (source0.tex lines 358--360, source label \texttt{None}), delivered together with the article's own complete original proof of it (source0.tex lines 361--404, one \texttt{proof} environment of 44 lines). No mathematics is written, repaired or re-derived: statement and proof are the article's own text line for line. Required context retained uncounted: mainthmSaeed (lines 168--176); mainthmSaeed2 (lines 301--309). Every definition, display and label of those lines is retained unchanged. Omitted from this package and not redistributed: 1--167, 177--300, 310--357, 405--525. Nothing inside a retained excerpt is omitted or altered: every retained body is delivered exactly as the article's lines are written, so no omission receipt of any kind accompanies this package. Citations: the retained excerpts carry no citation command at all, so no bibliography entry of the article is transcribed and no reference is redistributed. Reference adaptation: none was needed and none was made; every reference inside the delivered unit resolves inside it, so no unresolved reference or citation is produced. Printed slips kept literal: no printed word, symbol, hypothesis or number of the counted statement or of its proof was altered. Layout and class: the article's own class is \texttt{amsart} with packages including hyperref, amsmath, amssymb, latexsym, graphicx, fullpage, natbib, comment; the wrapper is the lane's pinned \texttt{amsart} editorial wrapper at 10pt on A4, which restates the theorem-like environments the retained text needs and adds the article's own macro definitions that the retained text needs (none), with extra wrapper packages (bm, bbm, dsfont, xspace, cancel, mathtools). The delivered numbering is the unit's own. Assets: no retained span contains a figure environment, a graphics-inclusion command, a PDF-page inclusion, a table or a TikZ picture; no image file is copied into this package, the package declares no asset dependency and no image file is redistributed with it. Licence: version arXiv:2101.04701v2 of this article is distributed under the Creative Commons Attribution 4.0 International licence, as recorded in the arXiv licence metadata for this exact version and shown on the abstract page https://arxiv.org/abs/2101.04701v2. A full rights notice is delivered at licenses/arxiv-2101.04701-CC-BY-4.0.txt. Characters: every retained excerpt is pure ASCII, so the delivered engine, XeLaTeX with its default Latin Modern text font, reports no missing character.
\begin{theorem}\label{mainthmSaeed}
	Let $r\geq 2$ be an integer, $C$ be a nonempty finite or infinite set, and
	$\tau: C \longrightarrow \{0,1,\ldots,r-1\}$ be an $r$-color-frequency  mapping.
	Then, for any $\tau$-admissible hypergraph ${\mathcal H}$ there exists a finite subset
	$B$ of $C$ for which
	$$\sum_{b\in B}\tau(b) \geq  |V({\mathcal H})|-{\rm alt}_r({\mathcal H}) .$$
	Besides, we have
	$$\chi_M(\tau,{\mathcal H})\geq \min \left\{|A|:\;   A \subseteq C, \mbox{ and } A \mbox{ is finite,}  \mbox{ and also } \sum_{a\in A}\tau(a) \geq  |V({\mathcal H})|-{\rm alt}_r({\mathcal H})\right\}.$$
\end{theorem}

\begin{theorem} \label{mainthmSaeed2}
	Let $C$ be a nonempty finite or infinite set, and $r$ be an integer with $r\geq 2$.
	Also, let the function $\tau: C \longrightarrow \{0,1,\ldots,r-1\}$ be an $r$-color-frequency mapping, and ${\mathcal H}$ be a $\tau$-admissible hypergraph. Then
	there exists a finite subset
	$B$ of $C$ for which
	$$\sum_{b\in B}\tau(b) \geq  {\rm ecd}^r({\mathcal H}) .$$
	Furthermore,
	$$\chi_M(\tau,{\mathcal H})\geq \min\left\{|A|:\   A \subseteq C, \mbox{ and } A \mbox{ is finite,}  \mbox{ and also } \sum_{a\in A}\tau(a) \geq  {\rm ecd}^r({\mathcal H})\right\}.$$
\end{theorem}

\begin{proposition}
	The lower bounds for $\chi_M(\tau,{\mathcal H})$ stated in Theorem \ref{mainthmSaeed} and Theorem \ref{mainthmSaeed2} are sharp.
\end{proposition}

\begin{proof}
	Let $n$, $k$, and $r$ be positive integers with $r\geq 2$ and
	$n\geq rk$. Also, let $\tau:\mathbb{N}\longrightarrow \{0,1,\ldots,r-1\}$
	be an $r$-color-frequency mapping such that $K_n^k$ is $\tau$-admissible and $r-1\in\tau(\mathbb{N})$.
	
	On account of Theorem \ref{mainthmSaeed}, the $\tau$-admissibility of $K_n^k$ implies
	$$\left\{|A|:\; A\subseteq \mathbb{N},\ A \mbox{ is finite, and } \sum_{a\in A}\tau(a) \geq  \left|V\left(K_n^k\right)\right|-{\rm alt}_r\left(K_n^k\right)\right\} \neq \varnothing ;$$
	and also,
	$$\chi_M\left(\tau,K_n^k\right) \geq \min \left\{|A|:\; A\subseteq \mathbb{N},\ A \mbox{ is finite, and } \sum_{a\in A}\tau(a) \geq  \left|V\left(K_n^k\right)\right|-{\rm alt}_r\left(K_n^k\right)\right\} .$$
	According to Theorem \ref{mainthmSaeed2}, the $\tau$-admissibility of $K_n^k$ indicates that
	$$\left\{|A|:\; A\subseteq \mathbb{N},\ A \mbox{ is finite, and } \sum_{a\in A}\tau(a) \geq  {\rm ecd}_r\left(K_n^k\right)\right\} \neq \varnothing ;$$
	and in addition,
	$$\chi_M\left(\tau,K_n^k\right) \geq \min \left\{|A|:\; A\subseteq \mathbb{N},\ A \mbox{ is finite, and } \sum_{a\in A}\tau(a) \geq  {\rm ecd}_r\left(K_n^k\right)\right\} .$$
	Now, since ${\rm alt}_r\left(K_n^k\right)= r(k-1)$ and ${\rm ecd}_r\left(K_n^k\right)= n - r(k-1)$,
	one finds that
	$$\left|V\left(K_n^k\right)\right|-{\rm alt}_r\left(K_n^k\right) = {\rm ecd}_r\left(K_n^k\right)= n - r(k-1) ;$$
	and therefore, we shall have established the Proposition if we show that
	$$\chi_M\left(\tau,K_n^k\right)\leq \min \left\{|A|:\; A\subseteq \mathbb{N},\ A \mbox{ is finite, and } \sum_{a\in A}\tau(a) \geq  n-r(k-1)\right\} .$$
	To do this, let us regard an arbitrary nonempty finite subset $A_0:=\{a_1,a_2,\ldots,a_t\}$ of $\mathbb{N}$ with the following four properties:
	\begin{itemize}
		\item $\displaystyle\sum_{a\in A_0}\tau(a) \geq  n-r(k-1) $; \\
		\item $\left| A_0 \right| = \min \left\{|A|:\; A\subseteq \mathbb{N},\ A \mbox{ is finite, and } \displaystyle\sum_{a\in A}\tau(a) \geq  n-r(k-1)\right\} $; \\
		\item $\tau(a_{1})\leq\tau(a_{2})\leq\cdots\leq\tau(a_{t})$; \\
		\item $\tau(a_{t})=r-1$ (Without loss of generality, the case 
		$\tau(a_{t})=r-1$
		may be considered. Otherwise, due to
		$r-1\in\tau(\mathbb{N})$, an element
		$b \in \mathbb{N}$ satisfying $\tau(b)=r-1$
		may be substituted for 
		$a_{t}$).
	\end{itemize}
	Now, let us consider $t$ pairwise disjoint sets $S_1, S_2, \ldots, S_t$ such that
	the following three conditions hold:
	\begin{itemize}
		\item $S_{1} \cup S_{2} \cup \dots \cup S_{t}=[n]$;
		\item For each $i$ in $\{1,2,\ldots ,t-1\}$ we have $|S_i|\leq \tau (a_i)$;
		\item $ |S_{t}| \leq \tau (a_t) + r(k-1) = r-1 + r(k-1)=rk-1$.
	\end{itemize}
	Also, let us define a mapping $c:\; E\left(K_n^k\right)\longrightarrow A_0$ such that each $e\in E\left(K_n^k\right)$ is mapped to $c(e):=a_{\varphi (e)}$ where
	$$ \varphi (e) := \min \left\{i:\; e\cap S_i\neq \varnothing \right\} . $$
	The mapping $c$ is an $\left(A_0 ,\tau \right)$-matching coloring of $K_n^k$; and therefore,
	$$\chi_M\left(\tau,K_n^k\right)\leq \left| A_0 \right| = \min \left\{|A|:\; A\subseteq \mathbb{N},\ A \mbox{ is finite, and } \sum_{a\in A}\tau(a) \geq  n-r(k-1)\right\} ;$$
	which is desired.
\end{proof}

\end{document}
